% ===== ЗАДАНИЕ, страница 1: общая шапка + УТВЕРЖДАЮ + заголовок =====
\thispagestyle{plain}
\begingroup
\singlespacing
\centering

Министерство науки и высшего образования РФ\\
Федеральное государственное автономное\\
образовательное учреждение высшего образования\\
\textbf{<<\varUniversity>>}\\[\baselineskip]
\varInstitute\\[\baselineskip]
<<\varDepartment>>

\vspace{2em}

\hfill
\begin{minipage}{0.46\textwidth}
\raggedright
УТВЕРЖДАЮ:\\
Заведующий кафедрой\\[0.6em]
\rule{2.8cm}{0.4pt}~~\varHead\\[0.6em]
<<\rule{0.7cm}{0.4pt}>>~\rule{2.4cm}{0.4pt}~\varYear~г.
\end{minipage}

\vspace{3\baselineskip}

\textbf{ЗАДАНИЕ}\\[0.3em]
\textbf{НА ВЫПУСКНУЮ КВАЛИФИКАЦИОННУЮ РАБОТУ}\\[0.3em]
\textbf{в форме бакалаврской работы}

\vfill
\endgroup
\clearpage

% ===== ЗАДАНИЕ, страница 2: форма «Студенту ...» =====
\thispagestyle{plain}
\begingroup
\singlespacing
\setlength{\parindent}{0pt}

Студенту\hspace{1em}\uline{\hspace{0.3em}\varStudentFull\hfill}\\[-0.3em]
\hspace*{\fill}{\fontsize{10}{12}\selectfont фамилия, имя, отчество}\hspace*{\fill}

\vspace{0.5em}
Группа\hspace{0.7em}\uline{~\varGroup~}\hspace{1.2em}Направление (специальность)\hspace{0.7em}\uline{~\varSpecCode~}\\[-0.3em]
\hspace*{1.6cm}{\fontsize{10}{12}\selectfont номер}\hspace*{\fill}{\fontsize{10}{12}\selectfont код}\hspace*{1.2cm}

\vspace{0.5em}
\uline{\varSpecName\hfill}\\[-0.3em]
\hspace*{\fill}{\fontsize{10}{12}\selectfont полное наименование}\hspace*{\fill}

\vspace{0.6em}
Тема выпускной квалификационной работы (далее~--- ВКР):\\
\uline{\varTheme\hfill}

\vspace{0.5em}
Утверждена приказом по университету\hspace{1em}№~\rule{1.2cm}{0.4pt}/с от~\rule{2.6cm}{0.4pt}

\vspace{0.5em}
Руководитель ВКР:\\
\uline{\varSupervisor~--- ст.\,преподаватель кафедры технологических машин и оборудования нефтегазового комплекса\hfill}\\[-0.3em]
\hspace*{\fill}{\fontsize{10}{12}\selectfont инициалы, фамилия, должность, ученое звание и место работы}\hspace*{\fill}

\vspace{0.6em}
\uline{1.~Теоретико-аналитический обзор предметной области исследования, проводимого в рамках ВКР, включая:\hfill}\\
\uline{--- описание конструкции и условий работы многоступенчатых центробежных насосов;\hfill}\\
\uline{--- описание гидродинамических подшипников скольжения и щелевых уплотнений рабочих колёс;\hfill}\\
\uline{--- анализ методов расчёта динамики роторов и постановку задачи исследования;\hfill}\\
\uline{2.~Разработка методики и численной модели динамики ротора;\hfill}\\
\uline{3.~Расчёт критических скоростей, орбит, перекоса шейки и анализ устойчивости ротора;\hfill}\\
\uline{4.~Нормативная оценка вибрационного состояния и анализ полученных результатов;\hfill}\\
\uline{Заключение по работе.\hfill}\\
\uline{Перечень графического материала: расчётная схема ротора; графики результатов моделирования (критические скорости, орбиты, логарифмический декремент); презентация к ВКР.\hfill}

\vspace{2.5em}
\begin{tabular}{@{}p{6.5cm}@{}>{\centering\arraybackslash}p{4.5cm}@{}>{\raggedleft\arraybackslash}p{\dimexpr\textwidth-11cm}@{}}
Руководитель ВКР & \rule{4.2cm}{0.4pt} & \varSupervisor \\[-0.3em]
 & {\fontsize{10}{12}\selectfont подпись} & \\[1.2cm]
Задание принял к исполнению & \rule{4.2cm}{0.4pt} & \varStudent \\[-0.3em]
 & {\fontsize{10}{12}\selectfont подпись студента} & \\
\end{tabular}

\vspace{1.5em}
\hspace*{\fill}<<\rule{0.8cm}{0.4pt}>>~\rule{3cm}{0.4pt}~\varYear~г.

\endgroup
\clearpage
